\section{Formulation dimension and convex contact geometry}\label{sec:foundations}

We study the number of auxiliary integer coordinates needed to retain a
nonlinear graph with a prescribed vertical error. This quantity measures
one resource of a formulation. It does not measure the time needed to
optimize over that formulation. In particular, an upper bound on integer
dimension need not give a small linear description, and a small linear
description need not have short rational coefficients. We state these
additional properties separately whenever they are available.

\subsection{The representation model}

Let $D\subset\R^n$ be compact and convex, and let $F:D\to\R^m$ be
continuous. Write $\Gamma_D(F)=\{(x,F(x)):x\in D\}$. A compact, convex,
centrally symmetric set $K\subset\R^m$ with nonempty interior is called
an \emph{error body}. The basic example is
$K=\prod_{j=1}^m[-\eps_j,\eps_j]$, with $\eps_j>0$.

\begin{definition}\label{def:dimension}
A graph relaxation with error body $K$ is a set $R\subset D\times\R^m$
satisfying
\[
 \Gamma_D(F)\subset R\subset
 \{(x,w):x\in D,\ w-F(x)\in K\}.
\]
A convex integer lift of integer dimension $p$ represents $R$ as
\begin{equation}\label{eq:convex-lift}
 R=\{(x,w):\exists a\in\R^q,\ z\in\Z^p,
                  \ (x,w,a,z)\in\mathcal C\},
 \qquad \mathcal C\text{ convex},
\end{equation}
where $q$ is finite but unrestricted. The convex set need not be closed,
its coefficients need not be rational, and the integer coordinates need
not be bounded. Let $\pconv(F,D;K)$ denote the least such $p$, or $+\infty$
if none exists. Let $\pbin(F,D;K)$ denote the least number of binary
coordinates when $\mathcal C$ is a polyhedron and $z\in\{0,1\}^p$.
For the common component tolerance $\eps$, abbreviate these quantities
as $\pconv(\eps)$ and $\pbin(\eps)$ when $F,D$ are fixed.
\end{definition}

One can equivalently allow inputs outside $D$ and measure error only on
$D$: intersection with $D\times\R^{m+q+p}$ preserves convexity. In the
binary linear model this equivalence applies when $D$ is a polytope,
which is the case for all linear constructions below. Always
$\pconv\le\pbin$. This notation deliberately reserves $\pbin$ for
\emph{linear} binary lifts; a lower bound valid for arbitrary convex
binary lifts will be stated explicitly when useful.

Adding an affine function to $F$ changes neither minimum: subtract it
from $w$ in the lift. Invertible affine changes of input coordinates also
preserve the minima when the domain is transformed along with the map.
For an invertible linear output map $A$, the pairs $(F,K)$ and $(AF,AK)$
are equivalent. Restricting a formulation to an affine input slice
preserves its integer dimension and gives the same error on that slice.
These observations justify all normalizations and restrictions below.

\begin{lemma}[Integer sections and parity contacts]\label{lem:parity}
For a lift \eqref{eq:convex-lift}, define its projected section at a
real index $z$ by
\[
 R_z=\{(x,w):\exists a,\ (x,w,a,z)\in\mathcal C\}.
\]
Every $R_z$ is convex, and
$(1-t)R_z+tR_{z'}\subset R_{(1-t)z+tz'}$ for $0\le t\le1$.
If the lift has error body $K$, there are at most $2^p$ compact sets
$S_\alpha$ covering $D$ such that, for every $x,y\in S_\alpha$,
\begin{equation}\label{eq:midpoint-error}
 \frac{F(x)+F(y)}2-F\left(\frac{x+y}2\right)\in K.
\end{equation}
More generally, if exact graph witnesses have integer indices $z_i$ and
nonnegative weights $\lambda_i$ summing to one satisfy
$\sum_i\lambda_i z_i\in\Z^p$, then
$\sum_i\lambda_iF(x_i)-F(\sum_i\lambda_i x_i)\in K$.
For binary lifts, fixing the binary vector gives at most $2^p$ convex
sections whose union is the relaxation.
\end{lemma}
\begin{proof}
Take witnesses for two projected points and average their continuous and
index coordinates. Convexity proves the section statements. For
$\alpha\in\{0,1\}^p$, let $T_\alpha$ contain those exact graph inputs
with some integer witness congruent to $\alpha$ modulo two. These sets
cover $D$; they may overlap and need not be measurable. The midpoint of
two witnesses in one class has an integer index, so its visible error
belongs to $K$. Set $S_\alpha=\overline{T_\alpha}$, taking closure in
$D$. Continuity of $F$ and closedness of $K$ preserve the inequality
for pairs of limits. These compact sets still cover $D$. The finite-mixture assertion follows
by taking the same convex combination of all witnesses.
\end{proof}

The parity mechanism is the midpoint obstruction of
\citet[Lemma 4.1]{lubin2022}; the proof above specifies the approximation
and closure conventions used here. A parity class is generally not a
convex section. Convexity of the common lift relates different integer
sections, whereas the parity argument extracts only the midpoint
condition \eqref{eq:midpoint-error}. Keeping this distinction is essential
when comparing unrestricted integers with binaries.

\begin{lemma}[Finite disjunction]\label{lem:disjunction}
A union of $N\ge1$ nonempty bounded polyhedra has a linear formulation
with $\lceil\log_2N\rceil$ binaries and unrestricted continuous size.
The same construction applies to finitely many polyhedra having a common
recession cone.
\end{lemma}
\begin{proof}
Write $Q_i=\{v:A_iv\le b_i\}$ and choose distinct codes
$c_i\in\{0,1\}^{\lceil\log_2N\rceil}$. Use continuous variables
$\lambda_i\ge0,v_i$, and impose
\[
 \sum_i\lambda_i=1,\quad A_iv_i\le\lambda_i b_i,\quad
 v=\sum_i v_i,\quad z=\sum_i\lambda_i c_i,\quad z\text{ binary}.
\]
An integral convex combination of binary vectors has every positive
weight on the same code. Thus there is one active member. If $Q_i$ is
bounded, $\lambda_i=0$ forces $v_i=0$. With a common recession cone,
the inactive $v_i$ instead belong to that cone and their sum can be
absorbed into the active member. The converse follows by placing weight
one on any selected member.
\end{proof}

\subsection{Squares, products, and shared input bits}

For a scalar function, write
$J_f(x,y)=(f(x)+f(y))/2-f((x+y)/2)$. If $f$ or $-f$ is
$\mu$-strongly convex, then
\begin{equation}\label{eq:strong-midpoint}
 |J_f(x,y)|\ge\frac\mu8\|x-y\|_2^2.
\end{equation}
Consequently a parity contact has diameter at most
$\sqrt{8\eps/\mu}$. Let $\omega_d$ denote the volume of the unit
Euclidean ball in $\R^d$. The isodiametric inequality and
\cref{lem:parity} yield, on a $d$-dimensional domain of volume $V$,
\begin{equation}\label{eq:strong-lower}
 \eps\ge\frac\mu2\left(\frac{V}{\omega_d}\right)^{2/d}2^{-2p/d}.
\end{equation}
The same estimate holds with $2^p$ replaced by the number of available
integer assignments when that number is finite. Strong convexity follows,
for example, from $\nabla^2f\succeq\mu I$ throughout the domain.

We will repeatedly use an elementary linear construction. For
$0\le u\le U$ and a binary $\beta$, the equality $v=\beta u$ is enforced by
\begin{equation}\label{eq:binary-product}
 0\le v\le U\beta,\qquad v\le u,\qquad
 v\ge u-U(1-\beta).
\end{equation}
For $0\le s\le h$ and $0\le t\le k$, use the four inequalities
\begin{equation}\label{eq:mccormick}
 q\ge0,\qquad q\ge ks+ht-hk,\qquad q\le ks,\qquad q\le ht
\end{equation}
These contain $q=st$ and give $|q-st|\le hk/4$.
Indeed after $s=ha,t=kb$, the upper error is
$hk\min\{a(1-b),b(1-a)\}\le hk/4$; the lower error follows by
replacing one of $a,b$ by its complement. This is the classical
McCormick relaxation \citep{mccormick1976}.
For a square on $[0,h]$, its specialization
\begin{equation}\label{eq:square-triangle}
 q\ge0,\qquad q\ge2hs-h^2,\qquad q\le hs
\end{equation}
contains the graph and has absolute error at most $h^2/4$.

To share precision across expressions, write
\begin{equation}\label{eq:prefix}
 y_i=A_i+\rho_i,\qquad
 A_i=\sum_{\ell=1}^{L_i}2^{-\ell}\beta_{i\ell},\qquad
 0\le\rho_i\le h_i=2^{-L_i}.
\end{equation}
Every $y_i\in[0,1]$ has a representation, including $1$, by using the
all-one prefix and residual $h_i$. Empty prefixes are allowed.
The identity
\begin{equation}\label{eq:product-split}
 y_i y_k=A_i y_k+\rho_i A_k+\rho_i\rho_k
\end{equation}
uses exact binary products in the first two terms and only one relaxed
residual product. It therefore has error at most $h_i h_k/4$, including
$i=k$ when \eqref{eq:square-triangle} is used. The same input bits and
approximate monomials can serve all outputs. A required monomial uses
$O(L_i+L_k+1)$ continuous variables and inequalities.

\begin{theorem}[Exact square law]\label{thm:square}
For $f(x)=x^2$ on $[0,1]$, every convex integer lift with $p$ integer
coordinates has error at least $2^{-2p-2}$. A binary linear lift attains
this error with $p$ binaries and $O(p+1)$ rows and variables. Thus, for
$\eps>0$,
\[
 \pconv(\eps)=\pbin(\eps)
 =\max\left\{0,\left\lceil\frac{\log_2(1/\eps)-2}{2}\right\rceil\right\}.
\]
\end{theorem}
\begin{proof}
Here $J_f(x,y)=(x-y)^2/4$. A compact set covering length at least
$2^{-p}$ has diameter at least $2^{-p}$, so the parity cover gives the
lower bound. For the upper bound, use \eqref{eq:prefix} in one variable
at depth $p$ and \eqref{eq:product-split} for its square. The only error
is the residual square error $h^2/4=2^{-2p-2}$. At a residual midpoint
its upper error equals this quantity, so the bound is attained. Exact
graph points lift by selecting their exact residual square values.
\end{proof}

This construction proves attainment directly. Shared binary discretization and residual McCormick relaxations also underlie the constructions of \citet{beach2024b}. The established sawtooth
formulation of \citet{beach2024} attains the same error with the same
binary count and linear size. Thus its integer count is also optimal
against arbitrary convex lifts with unrestricted integer ranges.

\begin{lemma}[Product contact geometry]\label{lem:product-width}
Suppose a compact $S\subset\R^2$ satisfies
$|(x_1-y_1)(x_2-y_2)|\le\delta$ for every $x,y\in S$.
If its coordinate widths are $d_1,d_2$, then
\[
 d_1d_2\le5\delta,\qquad \vol_2(S)\le4\delta\ln2.
\]
\end{lemma}
\begin{proof}
If $d_1=0$, both conclusions follow. Otherwise choose points $a,b$
attaining the extreme first coordinates. Their second coordinates differ
by at most $\delta/d_1$. Every other point is at first-coordinate
distance at least $d_1/2$ from one of them, hence its second coordinate
is within $2\delta/d_1$ of $a_2$ or $b_2$. This gives the width bound.
At first coordinate $a_1+t$, $0<t<d_1$, the vertical section lies in
two intervals, of lengths $2\delta/t$ and $2\delta/(d_1-t)$.
Integration gives
\[
 \vol_2(S)\le
 \int_0^{d_1}\min\left\{\frac{2\delta}{t},
                    \frac{2\delta}{d_1-t}\right\}\,dt
 =4\delta\ln2.
\]
\end{proof}

For the product $xy$ on $[0,1]^2$, $J_f(a,b)=(a_1-b_1)(a_2-b_2)/4$.
The lemma and parity therefore give
\begin{equation}\label{eq:product-finite}
 \frac{1}{16\ln2\,2^p}\le E_{
\min,\mathrm{conv}}(p)\le E_{\min,\mathrm{bin}}(p)\le2^{-p-2},
\end{equation}
where each $E_{\min}$ is the infimum achievable error with at most $p$
integer coordinates in its stated class. The upper bound uses depth
$p$ in the first input and depth zero in the second in
\eqref{eq:product-split}. Hence both minimum dimensions are
$\log_2(1/\eps)+O(1)$. For $0<\eps\le1/4$ the explicit bounds are
\[
 \left\lceil\log_2(1/\eps)-\log_2(16\ln2)\right\rceil
 \le\pconv(\eps)\le\pbin(\eps)
 \le\lceil\log_2(1/\eps)\rceil-2.
\]
Their difference is at most two. No exact optimum constant for the
whole product graph is claimed.

The width constant cannot be reduced from five to four. Put
$s=\sqrt5-2$ and consider
\[
 (0,(1-s)/2),\quad (1,(1+s)/2),\quad
 ((1+s)/2,0),\quad ((1-s)/2,1).
\]
Both widths are one, while all six pairwise difference products have
absolute value $s$, since $1-s^2=4s$. Any universal width constant must
therefore be at least $1/s=2+\sqrt5$. This obstruction concerns width;
it does not identify the best area constant or a realizable optimal
contact partition.

\begin{theorem}[Bilinear interaction graphs]\label{thm:graph-cover}
Let $G=(V,E)$ be a simple graph with $n$ vertices, and approximate all
outputs $x_ix_j$, $ij\in E$, on $[0,1]^V$ with positive tolerances
$\eps_{ij}$. For $C>0$ define
\[
 L_C=\min\left\{\sum_i s_i:s_i\ge0,\quad
 s_i+s_j\ge[\log_2(1/(C\eps_{ij}))]_+\quad(ij\in E)\right\},
\]
where $[t]_+=\max\{t,0\}$. If $s$ minimizes $L_4$, then
\begin{equation}\label{eq:graph-finite}
 \lceil L_{20}\rceil\le\pconv\le\pbin
 \le\sum_i\lceil s_i\rceil\le L_4+n.
\end{equation}
Let $\tau^*(G)=\min\{\sum_i a_i:a_i\ge0,\ a_i+a_j\ge1\ (ij\in E)\}$
be the fractional vertex-cover number. Then
$L_4-L_{20}\le\tau^*(G)\log_2 5$.
For common $\eps\downarrow0$, each minimum has asymptotic expansion
\[
 \tau^*(G)\log_2(1/\eps)+O_G(1).
\]
\end{theorem}
\begin{proof}
For any nonempty closed parity contact, let $d_i\le1$ be its coordinate
widths. Its two-coordinate projections satisfy
$d_id_j\le20\eps_{ij}$ by \cref{lem:product-width}. A zero width gives
zero $n$-volume. Otherwise $s_i=-\log_2d_i$ is feasible for $L_{20}$,
so the contact volume is at most $\prod_i d_i\le2^{-L_{20}}$.
Summing over at most $2^p$ contacts proves $p\ge L_{20}$.

For the upper bound give coordinate $i$ depth $L_i=\lceil s_i\rceil$
in \eqref{eq:prefix}, shared across all incident edges. Equation
\eqref{eq:product-split} has error at most
$2^{-L_i-L_j}/4\le\eps_{ij}$. The formulation has
$O(n+|E|+\sum_i\deg(i)L_i)$ continuous variables and rows. Adding
$\log_2(5)$ times an optimal fractional cover to an optimizer for
$L_{20}$ makes it feasible for $L_4$, because each edge demand increases
by at most $\log_2 5$. Finally, for $\eps<1/20$, common demands give
$L_C=\tau^*(G)\log_2(1/(C\eps))$, for $C=4,20$.
\end{proof}

For rational tolerances this is also an effective rational construction.
Replace each $L_4$ demand by
$d_{ij}=\max\{0,\lceil\log_2(1/(4\eps_{ij}))\rceil\}$,
computed by integer comparisons with powers of two, and solve the
resulting rational linear program. Rounding the demands increases its
optimum by at most $\tau^*(G)$. Thus a polynomial-time construction uses
at most $L_{20}+n+\tau^*(G)(1+\log_2 5)$ binaries and has polynomial
encoding length in the graph and tolerance data. The real logarithmic
linear program in \cref{thm:graph-cover} is an exact benchmark; computing
it exactly is unnecessary for this guarantee.

Stars have coefficient one; $K_{r,s}$ has coefficient $\min\{r,s\}$;
a matching of $k$ edges has coefficient $k$; an odd cycle on $2k+1$
vertices has coefficient $k+1/2$; and $K_n$, $n\ge2$, has coefficient
$n/2$. Thus shared bits and fractional allocation are both consequential.
On a general nondegenerate box replace $\eps_{ij}$ by
$\eps_{ij}/((u_i-l_i)(u_j-l_j))$. Affine terms introduced by the
normalization have no effect.

\section{Scalar quadratic rank and one-sided inertia}\label{sec:scalar-quadratic}

An indefinite quadratic is not strongly convex on its full domain, but
its graph still has the precision exponent of its rank. Cancellation
along some directions does not allow large contact volume in all
independent directions.

\begin{lemma}[Indefinite contact volume]\label{lem:indefinite-volume}
Let $M$ be a nonsingular real symmetric $d\times d$ matrix. If a compact
set $S\subset\R^d$ satisfies
$|\tfrac12(x-y)^TM(x-y)|\le\delta$ for all $x,y\in S$, then
\[
 \vol_d(S)\le
 \frac{2^d(3\sqrt d\,\delta)^{d/2}}{\sqrt{|\det M|}}.
\]
\end{lemma}
\begin{proof}
If the affine span has dimension below $d$, its volume is zero.
Otherwise choose $s_0,\ldots,s_d\in S$ maximizing
$|\det A|$, where $A=[s_1-s_0,\ldots,s_d-s_0]$. The maximum is
positive and attained. Replacing column $i$ by $s-s_0$ shows that each
coordinate of $A^{-1}(s-s_0)$ has absolute value at most one. Thus
$S\subset s_0+A[-1,1]^d$ and $\vol_d(S)\le2^d|\det A|$.
Writing $q(v)=v^TMv/2$, polarization gives for $G=A^TMA$
\[
 |G_{ii}|\le2\delta,\qquad
 |G_{ij}|=|q(s_i-s_0)+q(s_j-s_0)-q(s_i-s_j)|\le3\delta.
\]
Hadamard's inequality gives
$|\det G|\le(3\sqrt d\,\delta)^d$. Since
$\det G=(\det A)^2\det M$, the conclusion follows. This also excludes
a full-dimensional span when $\delta=0$.
\end{proof}

\begin{theorem}[Scalar quadratic rank]\label{thm:scalar-rank}
Let $B=\prod_i[l_i,u_i]$ be a bounded box with $l_i<u_i$, and
$f(x)=x^THx/2+a^Tx+b$, where $H=H^T$ has rank $r>0$.
Then each of $\pconv(\eps),\pbin(\eps)$ equals
\begin{equation}\label{eq:scalar-law}
 \frac r2\log_2(1/\eps)+O_{H,B}(1),\qquad\eps\downarrow0.
\end{equation}
A binary upper formulation uses $O_{H,B,n}(1+\log(1/\eps))$ rows and
variables. For any nonsingular principal submatrix $H[I,I]$ of order
$r$, every admissible $p$-integer lift satisfies
\[
 \eps\ge
 \frac{|\det H[I,I]|^{1/r}\bigl(\prod_{i\in I}(u_i-l_i)\bigr)^{2/r}}
      {48\sqrt r}\,2^{-2p/r}.
\]
If $r=0$, the exact graph is affine and both minima are zero.
\end{theorem}
\begin{proof}
A symmetric rank-$r$ matrix has a nonsingular principal submatrix of
order $r$: the sum of its principal minors of that order is the product
of its nonzero eigenvalues. Fix complementary coordinates in their
intervals. The remaining Hessian is $H[I,I]$ and the slice has volume
$V_I=\prod_{i\in I}(u_i-l_i)$. On any parity contact,
\[
 J_f(x,y)=\tfrac18(x-y)^TH[I,I](x-y),
\]
so \cref{lem:indefinite-volume} applies with $\delta=4\eps$.
The sum of at most $2^p$ contact volumes covers the slice. Rearranging
$V_I\le2^p2^r(12\sqrt r\,\eps)^{r/2}/\sqrt{|\det H[I,I]|}$
gives the explicit bound.

For the upper bound, diagonalize the quadratic and normalize its $r$
nonconstant linear coordinates to $y_j\in[0,1]$ on $B$. Then
$f=\ell+\sum_{j=1}^r c_jy_j^2$, where $\ell$ is affine and $c_j\ne0$.
Let $A=\sum_j|c_j|$. Apply \cref{thm:square} independently at common
depth $L=\max\{0,\lceil\tfrac12\log_2(A/(4\eps))\rceil\}$.
The output equation $w=\ell+\sum_jc_jt_j$ contains the graph and has
error at most $A2^{-2L-2}\le\eps$. It uses $rL$ binaries and
$O(r(L+1))$ auxiliary variables and rows, in addition to the affine
coordinate equations and original box. This proves the upper bound.
\end{proof}

For example $(x_1+\cdots+x_n)^2$ has rank one, independently of its
number of monomials. Conversely the inertia of a nonsingular indefinite
form does not change its whole-graph coefficient. Approximating a scalar
sum and approximating every monomial in that sum are different tasks.

\subsection{Epigraphs and hypographs}

An epigraph relaxation contains every $(x,w)$ with $x\in B,w\ge f(x)$
and admits only $w\ge f(x)-\eps$. Define
$p_{\mathrm{epi,conv}}(\eps)$ and $p_{\mathrm{epi,bin}}(\eps)$ with the
same lifting conventions. Hypograph relaxations contain $w\le f(x)$
and require $w\le f(x)+\eps$. The allowed output is unbounded in its
respective epigraph or hypograph direction.

We first record a continuous linear square underestimator of small
size. Put $G(t)=\min\{2t,2(1-t)\}$ on $[0,1]$, and
\[
 F_L(t)=t-\sum_{j=1}^L4^{-j}G^{\circ j}(t).
\]
Induction on the dyadic subdivision shows that $F_L$ is the linear
interpolant of $t^2$ at $k2^{-L}$: subdivision subtracts a triangular
function of height $4^{-L}$ at the new midpoints and zero at the old
knots. Therefore
\begin{equation}\label{eq:fold-square}
 0\le F_L(t)-t^2\le e_L:=4^{-L}/4.
\end{equation}
The epigraph $w\ge F_L(t)-e_L$ has a continuous linear lift of size
$O(L+1)$. Indeed replace the equalities $g_j=G(g_{j-1})$ by
\[
 g_0=t,\qquad 0\le g_j\le2g_{j-1},\qquad
 g_j\le2(1-g_{j-1}),\qquad
 w\ge t-\sum_{j=1}^L4^{-j}g_j-e_L.
\]
For fixed $t$, the weighted sum of $g_j$ is maximized by the exact
folding sequence. To check this, optimize successively backwards. If
$g_k$ decreases by $\Delta\ge0$ from its allowed maximum, the loss
$4^{-k}\Delta$ exceeds the largest possible gain in the optimized
future sum, whose Lipschitz constant is at most
$\sum_{j>k}4^{-j}2^{j-k}<4^{-k}$. This proves the claim by backward
induction; at the final level there is no future sum. Equation
\eqref{eq:fold-square} then gives a zero-binary epigraph relaxation of
error $e_L$. Using depth $L+1$ gives error $2^{-2L-4}$. This is the
folding principle underlying the sawtooth epigraph construction
\citep{beach2024}.

\begin{theorem}[One-sided inertia law]\label{thm:inertia}
For the scalar quadratic in \cref{thm:scalar-rank}, let
$k_-=n_-(H)$ and $k_+=n_+(H)$ be its negative and positive inertia.
If $k_->0$, both epigraph minima have expansion
\[
 \frac{k_-}{2}\log_2(1/\eps)+O_{H,B}(1).
\]
If $k_-=0$, both epigraph minima are zero for every $\eps>0$, and the
convex-lift minimum is zero also at exact accuracy. For hypographs,
replace $k_-$ by $k_+$. Binary linear upper formulations have
$O_{H,B,n}(1+\log(1/\eps))$ rows and variables.
\end{theorem}
\begin{proof}
Choose an orthonormal basis $T$ of the negative eigenspace and an
interior point $x_0$ of $B$. For a sufficiently small fixed $\rho>0$,
$x_0+T[-\rho,\rho]^{k_-}\subset B$. The restricted quadratic is
$\mu$-strongly concave for some $\mu>0$. Two graph witnesses in the
same parity class give an admitted midpoint below the true graph by
at least $\mu\|s-t\|^2/8$. The one-sided error bound thus yields
\eqref{eq:strong-lower} in dimension $k_-$ and proves the lower rate.

For the upper bound normalize a signed-square decomposition as
$f=\ell+\sum_{j\in P}c_jy_j^2-\sum_{j\in N}d_jy_j^2$,
where $c_j,d_j>0$ and $|N|=k_-$. Approximate each positive square
from below by the continuous linear epigraph just constructed, and
each negative square from above by the upper half of the graph
construction in \cref{thm:square}. The latter also contains the whole
square hypograph if the lower output bound is dropped. Impose
\[
 w\ge\ell+\sum_{j\in P}c_jt_j-\sum_{j\in N}d_jt_j.
\]
Depth $L=O(1+\log(1/\eps))$, chosen so that
$(\sum c_j+\sum d_j)2^{-2L-2}\le\eps$, gives the required error and
uses $k_-L$ binaries. Choosing $t_j=y_j^2$ lifts every exact epigraph
point, including arbitrarily large $w$. If $k_-=0$, no binaries are
used; the exact epigraph is itself convex. Apply the argument to $-f$
for the hypograph statement.
\end{proof}

For one product there is an exact one-sided convex-lift formula:
\begin{equation}\label{eq:product-one-sided}
 E_{\mathrm{epi,conv}}(p)=E_{\mathrm{hyp,conv}}(p)=2^{-2p-2}.
\end{equation}
To prove the lower bound restrict to $(x,y)=(t,1-t)$, where
$t(1-t)$ has concavity midpoint gap $(s-t)^2/4$; the interval parity
argument of \cref{thm:square} applies. For the upper bound use
$u=(x+y)/2$, $v=(x-y+1)/2$, so
$xy=u^2-v^2+v-1/4$. Represent the positive square epigraph exactly by
the convex inequality $t_+\ge u^2$, and use a depth-$p$ square
hypograph for $t_-$. The inequality
$w\ge t_+-t_-+v-1/4$ attains error $2^{-2p-2}$; take $u=1/2$ and
$v$ at a dyadic residual midpoint. Hypographs follow from the affine
transformation $(x,y,w)\mapsto(x,1-y,x-w)$. In a purely linear lift,
the exact positive square epigraph can be approached with arbitrarily
small extra error and no binaries. Exact finite-LP attainment at the
threshold in \eqref{eq:product-one-sided} is not asserted.

Zero binaries do not mean bounded continuous description size. Every
linear extended formulation with $M$ inequalities approximating the
square epigraph within $\eps$ satisfies
\begin{equation}\label{eq:epigraph-rows}
 M\ge\tfrac12\log_2(1/\eps)-1.
\end{equation}
To see this, write the projected lower boundary as the maximum of its
$q$ affine lower facets. On an active interval of length $h$, linearity
at its endpoints and the bounds $x^2-\eps\le g(x)\le x^2$ imply
$h^2/4\le\eps$. Hence $q\ge1/(2\sqrt\eps)$. Every projected facet
pulls back to a distinct exposed face of the extended polyhedron, which
has at most $2^M$ faces, since each face is determined by a subset of
tight inequalities. Equalities and lineality do not affect this count.
This proves \eqref{eq:epigraph-rows}; the folding construction gives
the matching logarithmic order.

\section{The exact precision exponent for quadratic systems}\label{sec:ncrank}

Let $F=(f_1,\ldots,f_m)$ with
\[
 f_j(x)=\tfrac12x^TH_jx+a_j^Tx+b_j,\qquad H_j=H_j^T,
\]
on a bounded full-dimensional box $B\subset\R^n$. The simultaneous
problem depends on the Hessian \emph{matrix space}, rather than on the
largest rank of a scalar combination of its members.

\subsection{Algebraic structure of a symmetric matrix space}

Let $\mathcal H=\operatorname{span}_{\C}\{H_1,\ldots,H_m\}$.
Its noncommutative rank $r$ is the rank of the pencil
$\sum_jt_jH_j$ over the complex free skew field, whose indeterminates
$t_j$ do not commute. This is a finite-dimensional invariant of the
coefficient matrices; the actual optimization inputs remain ordinary
real variables. We use three established characterizations
\citep[Theorems 1.4 and 1.17]{ggow2020}:
\begin{align}
 n-r&=\max_{U\subset\C^n}
       \bigl(\dim U-\dim\mathcal H(U)\bigr),
 &\mathcal H(U)&=\sum_jH_jU,\label{eq:shrunk}\\
 r=n&\ \Longleftrightarrow\ 
 \exists d,\ B_j\in\C^{d\times d}:
       \sum_j B_j\otimes H_j\text{ is invertible},\label{eq:evaluation}\\
 r=n&\ \Longleftrightarrow\ \kappa>0,
 &\kappa&=\inf_{P\succ0}
 \frac{\det\bigl(\sum_jH_jPH_j^*\bigr)}{\det P}.
 \label{eq:capacity}
\end{align}
The infimum in \eqref{eq:capacity} is over positive definite Hermitian
matrices. The involution on the free skew field conjugates scalars and
fixes the free variables, reversing the order of multiplication
\citep[Section 2.1]{volcic2021}. Thus a pencil with real symmetric
coefficients is Hermitian over that division ring. These algebraic
characterizations are imported results. The geometric and formulation
consequences are proved below.

\begin{lemma}[Principal compression]\label{lem:principal}
A Hermitian matrix of rank $r$ over a division ring with involution has
an invertible principal submatrix of order $r$. Consequently there is
a set $I$ of $r$ original coordinates such that
$\{H_j[I,I]\}_j$ has full noncommutative rank $r$.
\end{lemma}
\begin{proof}
Induct on matrix order. A nonzero diagonal entry is an invertible
one-by-one Hermitian pivot. If all diagonal entries vanish but the matrix
is nonzero, a nonzero off-diagonal entry $a$ supplies the principal pivot
\[
 D=\begin{pmatrix}0&a\\a^*&0\end{pmatrix},\qquad
 D^{-1}=\begin{pmatrix}0&(a^*)^{-1}\\a^{-1}&0\end{pmatrix}.
\]
After placing a pivot first write the matrix as
$\left(\begin{smallmatrix}D&E\\E^*&G\end{smallmatrix}\right)$.
Its Schur complement $S=G-E^*D^{-1}E$ is Hermitian and block elimination
gives $\rank M=\operatorname{order}(D)+\rank S$. Induction selects a
principal invertible submatrix $S[J,J]$ of order $\rank S$.
The original principal submatrix on the pivot indices and $J$ has this
Schur complement, so it is invertible. The zero matrix uses the empty
submatrix. Apply the lemma to the free-field pencil to obtain $I$.
\end{proof}

\begin{lemma}[Real symmetric shrinking]\label{lem:symmetric-shrink}
There are mutually orthogonal real subspaces $Z,W,Q$ with
$\R^n=Z\oplus W\oplus Q$ such that
\begin{equation}\label{eq:shrinking-blocks}
 H_j Z\subset W\quad\text{for every }j,
 \qquad 2\dim W+\dim Q=r.
\end{equation}
In orthonormal coordinates ordered as $Z,W,Q$, each Hessian has form
\begin{equation}\label{eq:block-pattern}
 \begin{pmatrix}
 0&B_j&0\\ B_j^T&C_j&D_j\\0&D_j^T&E_j
 \end{pmatrix}.
\end{equation}
\end{lemma}
\begin{proof}
Put $d=n-r$ and $g(U)=\dim U-\dim\mathcal H(U)$. The identities
$\mathcal H(U+V)=\mathcal H(U)+\mathcal H(V)$ and
$\mathcal H(U\cap V)\subset\mathcal H(U)\cap\mathcal H(V)$ show that
$g(U+V)+g(U\cap V)\ge g(U)+g(V)$. A maximizing complex subspace exists
because the deficiencies take finitely many integer values. Since the
matrices are real, its conjugate also has deficiency $d$.
Supermodularity forces both the sum and intersection with its conjugate
to have deficiency $d$. The sum is conjugation invariant and hence is
the complexification of a real space $U_0$: for any member, its real
and imaginary parts belong to the sum and span it over $\C$.
The image likewise complexifies its real image. Consequently a real
pair $U,V=\sum_jH_jU$ has $\dim U-\dim V=d$.

Set $Z=U\cap V^\perp$ and $W=V\cap U^\perp$. The restricted orthogonal
projections $P_V:U\to V$ and $P_U:V\to U$ are adjoints, and therefore
have the same rank. Their kernels are $Z,W$, so
$\dim Z-\dim W=d$. For $z\in Z$ we have $H_jz\in V$, and for $u\in U$
\[
 u^TH_jz=(H_ju)^Tz=0.
\]
Thus $H_jz\in V\cap U^\perp=W$. Moreover $Z\perp W$.
Let $Q=(Z+W)^\perp$. The dimension identity gives
$2\dim W+\dim Q=n-d=r$, and symmetry gives the displayed zero blocks.
\end{proof}

Assign an exponent $\alpha_i$ equal to zero on $Z$, one on $W$, and
one half on $Q$. Every nonzero quadratic coefficient has
$\alpha_i+\alpha_k\ge1$, including square coefficients with $i=k$,
while
\begin{equation}\label{eq:exponent-sum}
 \sum_i\alpha_i=r/2.
\end{equation}
This converts an algebraic rank witness into a precision allocation.

\subsection{Covariance bounds for contact volume}

\begin{lemma}[Contact covariance]\label{lem:covariance}
Let $S\subset\R^s$ be compact with positive volume, and let $\Sigma$
be the covariance matrix of its uniform distribution. For real symmetric
$G_j$, suppose
$|\tfrac12(x-y)^TG_j(x-y)|\le\delta_j$ for all $x,y\in S$.
Then $\Sigma\succ0$ and
\begin{align}
 \sum_j\tr(G_j\Sigma G_j\Sigma)&\le\sum_j\delta_j^2,
 \label{eq:covariance-energy}\\
 \vol_s(S)&\le\omega_s(s+2)^{s/2}\sqrt{\det\Sigma}.
 \label{eq:covariance-volume}
\end{align}
\end{lemma}
\begin{proof}
A singular covariance would place almost every point in an affine
hyperplane, contradicting positive volume. Center independent uniform
points $X,Y$ at zero. Expanding the square and using independence gives
\[
 \mathbb E\left[\left(\tfrac12(X-Y)^TG_j(X-Y)\right)^2\right]
 =\tfrac12\mathbb E[(X^TG_jX)^2]
  +\tfrac12\bigl(\mathbb E[X^TG_jX]\bigr)^2
  +\tr(G_j\Sigma G_j\Sigma).
\]
Mixed terms containing $X^TG_jY$ vanish after conditioning on one
centered factor. This proves \eqref{eq:covariance-energy}.
For \eqref{eq:covariance-volume}, whiten the centered set by
$\Sigma^{-1/2}$. Its mean squared norm is $s$. Among sets of a fixed
volume, a centered ball minimizes the integral of squared norm:
subtracting the ball's integral pairs exterior points with missing
interior points, and the former have no smaller squared norm. A ball
of radius $R$ has mean squared norm $sR^2/(s+2)$. Thus the whitened
set has volume at most $\omega_s(s+2)^{s/2}$, proving the result.
\end{proof}

Suppose the tuple $G_j$ has full noncommutative rank $s$. Write
$T(P)=\sum_jG_jPG_j$ and let $\kappa>0$ be its capacity. The ordinary
eigenvalue arithmetic--geometric mean inequality implies
\begin{equation}\label{eq:energy-det}
 \tr(\Sigma T(\Sigma))
 \ge s[\det\Sigma\det T(\Sigma)]^{1/s}
 \ge s\kappa^{1/s}(\det\Sigma)^{2/s}.
\end{equation}
Here $T(\Sigma)$ is positive definite by positive capacity, and the
first inequality is applied to $\Sigma^{1/2}T(\Sigma)\Sigma^{1/2}$.
Combining \cref{lem:covariance} with $\delta_j=4\eps$ yields
\begin{equation}\label{eq:ncrank-volume}
 \vol_s(S)\le C_s\eps^{s/2},\qquad
 C_s=\omega_s(s+2)^{s/2}
       \left(\frac{16m}{s\kappa^{1/s}}\right)^{s/4}.
\end{equation}
Zero-volume sets satisfy this inequality as well.

There is also an elementary way to obtain a positive constant in
\eqref{eq:energy-det}, without using capacity. For an orthogonal $O$
define $M_{ab}(O)=\sum_j((O^TG_jO)_{ab})^2$. If the support bipartite
graph of $M(O)$ failed Hall's condition, a coordinate column subspace
would be mapped by every $O^TG_jO$ into a smaller row subspace.
This would be a real shrunk subspace, contradicting \eqref{eq:shrunk}.
Hence its permanent is positive. Continuity and compactness of the
orthogonal group give
$\mu=\min_O\operatorname{per}M(O)>0$.
If $\Sigma=O\diag(\lambda)O^T$, some permutation $\pi$ has
$\prod_aM_{a,\pi(a)}\ge\mu/s!$. Therefore
\[
 \tr(\Sigma T(\Sigma))
 =\sum_{a,b}M_{ab}\lambda_a\lambda_b
 \ge s(\mu/s!)^{1/s}(\det\Sigma)^{2/s}.
\]
Thus the qualitative rate uses only the shrinking characterization.
The constant $\mu/s!$ need not equal capacity; capacity remains useful
for explicit rational lower estimates.

\begin{theorem}[Quadratic-system precision law]\label{thm:ncrank}
For a fixed real quadratic system $F$ on a full-dimensional bounded
box $B$, let $r$ be the noncommutative rank of its Hessian space.
If $r>0$, each minimum has expansion
\begin{equation}\label{eq:ncrank-law}
 \pconv(\eps),\ \pbin(\eps)
   =\frac r2\log_2(1/\eps)+O_{F,B}(1).
\end{equation}
The binary upper formulation has
$O_{F,B,n,m}(1+\log(1/\eps))$ rows and variables, with real coefficients.
If $r=0$, all outputs are affine and both minima are exactly zero.
For a principal compression $I$ from \cref{lem:principal}, write
$V_I=\prod_{i\in I}(u_i-l_i)$ and let $\kappa_I>0$ be the capacity of
the compressed tuple. Every admissible $p$-integer lift satisfies
\begin{equation}\label{eq:ncrank-finite}
 \eps\ge
 \frac{\sqrt{r\kappa_I^{1/r}/m}\ V_I^{2/r}}
      {4(r+2)\omega_r^{2/r}}\,2^{-2p/r}.
\end{equation}
\end{theorem}
\begin{proof}
Fix the coordinates outside $I$ in their intervals. The restricted
quadratic system has Hessians $G_j=H_j[I,I]$ of full noncommutative
rank $r$. On a parity contact its midpoint errors give
$|\tfrac12(x-y)^TG_j(x-y)|\le4\eps$. Equations
\eqref{eq:ncrank-volume} and \cref{lem:parity} imply
$V_I\le2^pC_r\eps^{r/2}$. Rearranging proves
\eqref{eq:ncrank-finite} and the lower rate.

For the upper bound take the coordinates of
\cref{lem:symmetric-shrink}, enclose the transformed box in a box,
and normalize its coordinates to $[0,1]^n$. Retain the original box
inequalities and coordinate identities. Diagonal rescaling and
translation do not create quadratic coefficients in zero blocks.
Write each output as an affine function plus
$\sum_{i\le k}c_{j,ik}y_iy_k$, and put
$C=\max_j\sum_{i\le k}|c_{j,ik}|>0$.
Take $T=\max\{0,\log_2(C/(4\eps))\}$ and depths
$L_i=\lceil\alpha_iT\rceil$ in \eqref{eq:prefix}. For every nonzero
monomial $2^{-L_i-L_k}\le2^{-T}$, by \eqref{eq:exponent-sum} and its
preceding coefficient condition. Shared residual-product relaxations
therefore give every output error at most $C2^{-T}/4\le\eps$.
Exact graph points lift with exact residual values. The binary count is
\[
 \sum_iL_i\le\tfrac r2T+n,
\]
and at most $n(n+1)/2$ monomials are needed. The size statement follows
from \eqref{eq:product-split}.
\end{proof}

Equations such as \eqref{eq:ncrank-law} assert the same asymptotic
expansion for two minima; they do not assert equality at every accuracy.
Also, a graph on a lower-dimensional feasible set need not obey the
full-box lower bound. In particular the theorem cannot be transferred
to the zero set of the output equations without a separate argument.

\subsection{A certificate and a strict scalar-rank gap}

If $\sum_jH_j^2=cI_n$ with $c>0$, one can replace capacity by an explicit
constant. In a covariance eigenbasis set
$M_{ab}=\sum_j((O^TH_jO)_{ab})^2$. Its row and column sums are $c$.
Weighted arithmetic--geometric mean with weights $M_{ab}/(cn)$ gives
\[
 \sum_j\tr(H_j\Sigma H_j\Sigma)
 \ge cn(\det\Sigma)^{2/n}.
\]
Thus \cref{lem:covariance} gives the finite bound
\begin{equation}\label{eq:sos-certificate}
 \eps\ge
 \frac{\sqrt{cn/m}\,\vol_n(B)^{2/n}}
      {4(n+2)\omega_n^{2/n}}\,2^{-2p/n}.
\end{equation}
Together with uniform half-depths in all coordinates, this proves
coefficient $n/2$ directly. It also certifies full noncommutative rank:
for a real or complex subspace $U$ with image contained in $V$, writing
$P_U,P_V$ for orthogonal projections gives
\[
 c\dim U=\sum_j\|H_jP_U\|_F^2
 \le\sum_j\|P_VH_j\|_F^2=c\dim V,
\]
where $P_VH_jP_U=H_jP_U$ and symmetry is used in the second total.
There is no shrunk subspace.

For $F(x,y)=x\times y$, $x,y\in[0,1]^3$, let
\[
 A(a)=\begin{pmatrix}0&a_3&-a_2\\-a_3&0&a_1\\a_2&-a_1&0\end{pmatrix},
 \qquad H(a)=\begin{pmatrix}0&A(a)\\-A(a)&0\end{pmatrix}.
\]
Then $a^TF=x^TA(a)y$ has Hessian $H(a)$.
Since $A(a)^TA(a)=\|a\|^2I-aa^T$, every nonzero scalar combination
has rank four. Yet $\sum_{j=1}^3H(e_j)^2=2I_6$, so both simultaneous
minimum dimensions have expansion $3\log_2(1/\eps)+O(1)$.
The explicit lower bound is
$\eps\ge(16\omega_6^{1/3})^{-1}2^{-p/3}$, where $\omega_6=\pi^3/6$.
For $k$ independent copies the coefficient is $3k$, while half the
maximum scalar rank is $2k$. Thus maximum scalar rank misses an
arbitrarily large additive amount in the precision coefficient.
The alternating matrix space itself is classical; the distinction here
is between its simultaneous graph and any single scalar output.

The two independently proved laws also identify the algebraic invariant for an interaction graph: comparing \cref{thm:graph-cover,thm:ncrank} as $\eps\downarrow0$ gives $\ncrank\operatorname{span}_{ij\in E}\{e_ie_j^T+e_je_i^T\}=2\tau^*(G)$. This identity follows by equality of the precision coefficients and needs no additional graph-rank theorem.

\subsection{Rational construction and its scope}\label{subsec:rational-scope}

For rational data there is an effective coordinate transformation.
The algorithm for noncommutative rank due to
\citet[Theorem 1.5 and Lemma 5.3]{iqs2018} returns a maximal rational
shrunk subspace with polynomial bit complexity and preserves its rank
under field extension. Rational bases of $Z,W,Q$ in
\cref{lem:symmetric-shrink} suffice: they need not be orthonormal,
because their spaces are orthogonal. Rational nullspace computations
therefore preserve \eqref{eq:block-pattern} without introducing square
roots.

More precisely, for a rational quadratic system and rational
full-dimensional box with total binary encoding length $s$, there
are universal polynomials $P,Q$ such that for every integer $k\ge0$
\begin{equation}\label{eq:rational-rate-preview}
 \tfrac r2 k-Q(s)\le\pconv(2^{-k})\le\pbin(2^{-k})
                    \le\tfrac r2 k+P(s).
\end{equation}
A deterministic algorithm constructs the upper relaxation in time and
total encoding length polynomial in $s+k$. Its individual coefficient
lengths are $\operatorname{poly}(s)+O(k)$, and its number of rows and
variables is $\operatorname{poly}(s)(k+\operatorname{poly}(s)+1)$.
The detailed bit bounds, including a uniform lower estimate from the
integral capacity bound \citep[Theorem 2.18]{ggow2020}, are developed
with the finite-accuracy quadratic constructions later in the paper.
This statement concerns construction, not polynomial-time optimization
of the resulting MILP. For arbitrary real data
\cref{thm:ncrank} makes no claim about finding or encoding the coordinate
change in a bit model.

\section{Smooth maps, varying null spaces, and perspectives}\label{sec:smooth}

For a smooth nonlinear map, Hessians at different inputs need not have
a common null space or a common optimal precision allocation. We
therefore distinguish two ranks. For $F=(f_1,\ldots,f_m)$ smooth on a
neighborhood of a full-dimensional compact box $B$, set
\begin{align*}
 r_{\mathrm{loc}}&=\max_{x\in\operatorname{int}B}
       \ncrank\operatorname{span}_j\{\nabla^2 f_j(x)\},\\
 r_{\mathrm{all}}&=\ncrank\operatorname{span}_{x\in B,j}
                                      \{\nabla^2 f_j(x)\}.
\end{align*}
The global span is finite-dimensional, being a space of symmetric
$n\times n$ matrices. These are matrix-space ranks, not dimensions of
the spans, and $0\le r_{\mathrm{loc}}\le r_{\mathrm{all}}\le n$.

\subsection{A local contact-volume estimate}

We give the analytic ingredient with its local proof, since it is what
allows contacts of arbitrary aspect ratio. The estimate is the classical
nondegenerate mixed-Hessian oscillatory-integral bound
\citep[Theorem 1.1]{hormander1973}; see also
\citet[Theorem A, pp.~50--51]{wolff2002}.

\begin{lemma}[Local oscillatory estimate]\label{lem:oscillatory}
Suppose a real $C^\infty$ phase $\Phi(X,Y)$, with $X,Y\in\R^N$, has
invertible mixed Hessian at $(X_0,Y_0)$. On a sufficiently small product
neighborhood, for every smooth compactly supported amplitude $a$,
\[
 T_\lambda g(X)=\int e^{i\lambda\Phi(X,Y)}a(X,Y)g(Y)\,dY
 \quad\text{satisfies}\quad
 \|T_\lambda\|_{L^2\to L^2}\le C\lambda^{-N/2},\qquad\lambda\ge1.
\]
The constant depends on the fixed phase, amplitude, and neighborhood.
\end{lemma}
\begin{proof}
Let $A_0=\Phi_{XY}(X_0,Y_0)$ and let $\sigma>0$ be its least singular
value. Shrink a convex product neighborhood until
$\|\Phi_{XY}-A_0\|\le\sigma/2$. The gradient in $Y$ of
$\psi(Y)=\Phi(X,Y)-\Phi(Z,Y)$ is the integral of
$\Phi_{XY}(Z+t(X-Z),Y)^T(X-Z)$, hence has norm at least
$(\sigma/2)|X-Z|$. Every fixed-order $Y$ derivative of $\psi$ is
$O(|X-Z|)$, by smoothness and integration in $X$.
The kernel of $T_\lambda T_\lambda^*$ is the integral of
$e^{i\lambda\psi(Y)}a(X,Y)\overline{a(Z,Y)}$. For $X\ne Z$ integrate
by parts with
\[
 \mathcal L=\frac{\nabla\psi\cdot\nabla_Y}
                  {i\lambda|\nabla\psi|^2},
 \qquad \mathcal L e^{i\lambda\psi}=e^{i\lambda\psi}.
\]
The derivative bounds show that each integration gains a factor
$O((\lambda|X-Z|)^{-1})$; the amplitude has compact support, so there
are no boundary terms. Combining any fixed number $M$ of integrations
with the trivial kernel bound gives
$|K_\lambda(X,Z)|\le C_M(1+\lambda|X-Z|)^{-M}$.
Choose $M>N$. Both Schur integrals are $O(\lambda^{-N})$, so Schur's
test bounds $\|T_\lambda T_\lambda^*\|$ by that quantity. Taking the
square root proves the result. Indicator functions are permitted by
$L^2$ continuity; no boundary regularity of their support is needed.
\end{proof}

\begin{lemma}[Smooth contact volume]\label{lem:smooth-volume}
If the Hessian tuple of $F$ has full noncommutative rank $n$ at an
interior point $x_0$, then there are a compact box $Q$ about $x_0$ and
constants $C,\eps_0>0$ such that every compact $S\subset Q$ with
\[
 |f_j(x)+f_j(y)-2f_j((x+y)/2)|\le2\eps
 \quad(x,y\in S,\ 1\le j\le m)
\]
has volume at most $C\eps^{n/2}$ for $0<\eps\le\eps_0$.
\end{lemma}
\begin{proof}
By \eqref{eq:evaluation} there exist $d$ and matrices $B_j$ for which
$\sum_j B_j\otimes\nabla^2f_j(x_0)$ is invertible. They can be chosen
real: its determinant is a real-coefficient polynomial in the matrix
entries, and a nonzero real polynomial cannot vanish at every real
tuple. Put $D_j(x,y)=f_j(x)+f_j(y)-2f_j((x+y)/2)$, and for
$X=(x_1,\ldots,x_d)$, $Y=(y_1,\ldots,y_d)$ define
\[
 \Phi(X,Y)=\sum_{a,b,j}(B_j)_{ab}D_j(x_a,y_b).
\]
At the repeated base point its mixed Hessian is
$-\tfrac12\sum_jB_j\otimes\nabla^2f_j(x_0)$.
Choose $Q$ small enough that a smooth cutoff equal to one on
$Q^d\times Q^d$ has support in the neighborhood of
\cref{lem:oscillatory}. Set $K_0=\sum_{a,b,j}|(B_j)_{ab}|>0$.
On $S^d\times S^d$, $|\Phi|\le2K_0\eps$, so with
$\lambda=(4K_0\eps)^{-1}\ge1$ the real part of the exponential is
at least $1/2$. Writing $v=\vol_n(S)$ gives
\[
 \tfrac12v^{2d}
 \le\left|\langle 1_{S^d},T_\lambda1_{S^d}\rangle\right|
 \le C\lambda^{-nd/2}v^d
 =C(4K_0\eps)^{nd/2}v^d.
\]
If $v>0$, divide by $v^d$ and take the $d$th root. If $v=0$ the
conclusion is immediate. The evaluation dimension cancels from the
exponent, although it affects the fixed constant.
\end{proof}

A direct small perturbation of the quadratic determinant argument would
not supply this lemma: a remainder controlled by
$\eta\|x-y\|^2$ can overwhelm a determinant bound on highly anisotropic
sets. The oscillatory estimate controls arbitrary compact contacts
without an aspect-ratio assumption.

\begin{theorem}[Local and global smooth ranks]\label{thm:smooth-ranks}
Suppose $F$ is real $C^\infty$ on an open neighborhood of $B$. Then
\begin{equation}\label{eq:smooth-sandwich}
 \frac{r_{\mathrm{loc}}}{2}\log_2(1/\eps)-O_{F,B}(1)
 \le\pconv(\eps)\le\pbin(\eps)
 \le\frac{r_{\mathrm{all}}}{2}\log_2(1/\eps)+O_{F,B}(1).
\end{equation}
If both ranks are $r$, both minima have leading coefficient $r/2$.
In particular, full pointwise rank at one interior point gives
coefficient $n/2$. The general upper construction may have size
polynomial in $1/\eps$ for fixed data. If every output is a polynomial
of degree at most a fixed $D\ge2$, its upper formulation instead has
$O_{F,B,n,m,D}((1+\log(1/\eps))^D)$ rows and variables, with the same
binary count. Affine maps have exact zero-binary linear formulations.
\end{theorem}
\begin{proof}
Choose an interior point of maximum local rank, which exists because
occurring ranks are integers in a finite range. For positive rank,
\cref{lem:principal} supplies a principal slice of dimension
$r_{\mathrm{loc}}$ whose Hessian tuple has full rank at that point.
Apply \cref{lem:smooth-volume} on a smaller box in the slice and use
the compact parity cover. Its positive volume is at most
$2^pC\eps^{r_{\mathrm{loc}}/2}$, proving the lower bound. Rank zero
requires only $p\ge0$.

For the upper bound apply \cref{lem:symmetric-shrink} to a finite
basis of the global Hessian span. Its coordinates give fixed exponents
$\alpha_i\in\{0,1,1/2\}$ summing to $r_{\mathrm{all}}/2$, and every
possibly nonzero Hessian entry has endpoint exponent sum at least one.
Enclose the transformed domain $P$ in a box of side lengths $b_i>0$.
Divide side $i$ into $2^{\lceil\alpha_iT\rceil}$ intervals and intersect
the cells with $P$. For each nonempty cell choose a center $c$ in that
intersection. Along the segment from $c$ to any point $x$ of the cell,
every Hessian is bounded and has the stated zero blocks. The integral
Taylor remainder therefore satisfies, uniformly over cells and outputs,
\[
 |f_j(x)-f_j(c)-\nabla f_j(c)^T(x-c)|\le C_0,2^{-T}.
\]
Indeed each contributing product $(x-c)_i(x-c)_k$ has absolute value
at most $b_i b_k2^{-T}$. Taking $2C_0,2^{-T}\le\eps$, use an affine
output band of radius $\eps/2$ on each cell. Each band contains its
exact graph and has error at most $\eps$. There are at most
$2^{\sum_i\lceil\alpha_iT\rceil}$ bounded polyhedral members, so
\cref{lem:disjunction} proves the upper binary count. If all Hessians
vanish, the map is affine on the convex box.

For polynomial outputs, each forbidden Hessian entry vanishes on an
open set and thus identically. The same structure is valid on the
whole enclosing box, which may now be normalized to $[0,1]^n$.
Use the input prefixes \eqref{eq:prefix}, with
$L_i=\lceil\alpha_iT\rceil$, and the first-order expression
\begin{equation}\label{eq:polynomial-prefix}
 P_j(A,\rho)=f_j(A)+\sum_i\partial_i f_j(A)\rho_i.
\end{equation}
The Taylor error is again at most $C_0,2^{-T}$, so impose
$|w_j-P_j(A,\rho)|\le\eps/2$. Expand \eqref{eq:polynomial-prefix}.
Each term is a product of at most $D$ bits, or at most $D-1$ bits and
one residual. Repeated bits can be removed. For a set $J$ of distinct
bits, continuous variables satisfying
\[
 0\le v_J\le1,\quad v_J\le\beta_b\ (b\in J),\quad
 v_J\ge\sum_{b\in J}\beta_b-|J|+1
\]
are forced to equal the bit product whenever the input bits are integral.
A product $v_J\rho_i$ is then represented by
\eqref{eq:binary-product}; $v_J$ needs no separate integrality declaration.
Empty products equal one. With fixed $D,n,m$, the expansion contains
$O((1+\sum_iL_i)^D)$ terms and $O(D)$ rows per term. All graph points
are retained because the affine band covers the Taylor remainder.
No additional integer coordinates are introduced.
\end{proof}

A product of $k\ge2$ positive inputs has coefficient $k/2$. At a point
with all coordinates equal to $a>0$, its Hessian has zero diagonal and
off-diagonal entries $a^{k-2}$, with eigenvalues
$(k-1)a^{k-2}$ once and $-a^{k-2}$ with multiplicity $k-1$.
A rank-deficient example is
\[
 f(z_1,z_2,w,t)=z_1w^3+z_2(w+w^2)+wt^2+t^4.
\]
Its Hessians have $Z=(z_1,z_2)$, $W=w$, $Q=t$ in
\eqref{eq:block-pattern}, so their global rank is at most three by the
shrinking characterization. At $(1,1,1,1)$ the Hessian is
\[
 \begin{pmatrix}0&0&3&0\\0&0&3&0\\3&3&8&2\\0&0&2&14\end{pmatrix},
\]
which has rank three. Thus its coefficient is $3/2$ on a box containing
that point in its interior, with an upper formulation of size
$O((1+\log(1/\eps))^4)$.

There is a representation-dependent computational boundary. Given a
rational arithmetic circuit for a polynomial $g(z)$ in $d$ variables,
form
\[
 F(x,y,z)=(xyg(z),z_1^2,\ldots,z_d^2)
 \quad\text{on }[-1,1]^{d+2}.
\]
If $g\equiv0$, its precision coefficient is $d/2$. Otherwise
$g(z_0)\ne0$ at some interior point, and at $(0,0,z_0)$ the Hessian of
the scalar combination $xyg(z)+\sum_i z_i^2$ is
\[
 \begin{pmatrix}0&g(z_0)&0\\g(z_0)&0&0\\0&0&2I_d\end{pmatrix},
\]
which is invertible. The coefficient is then $(d+2)/2$ by
\cref{thm:smooth-ranks}. Computing the exact coefficient, or an additive
approximation with error strictly below $1/2$, for such circuits would
therefore decide polynomial identity testing. This is a reduction, not
an unconditional running-time lower bound. In a dense coefficient
representation identity testing itself is immediate.

\subsection{Constant scalar Hessian rank}

The global-span upper bound in \cref{thm:smooth-ranks} need not be sharp.
For scalar maps with constant Hessian rank one can follow the local
null spaces by polyhedral tubes. The affine-gradient-fiber geometry is
classical \citep[Definition 1.1 and the following paragraph]{nicola2008};
we derive it locally to show why these tubes are linear objects.

\begin{theorem}[Constant-rank scalar law]\label{thm:constant-rank}
Let $f$ be $C^\infty$ on a neighborhood $U$ of the compact
full-dimensional box $B$, and assume $\rank\nabla^2f(x)=r$ for every
$x\in U$. Both minimum dimensions have expansion
\[
 \tfrac r2\log_2(1/\eps)+O_{f,B}(1).
\]
The upper bound is a union of $O_{f,B}(\eps^{-r/2})$ bounded polyhedra,
encoded with the logarithm of that many binary codes. It does not
assert a formulation of size polynomial in $\log(1/\eps)$.
\end{theorem}
\begin{proof}
Rank zero gives an affine function on $B$. For $r>0$ the lower bound
is the scalar case of \cref{thm:smooth-ranks}. At any point of $B$
choose coordinates $(u,v)\in\R^r\times\R^{n-r}$ with $f_{uu}$
invertible. The inverse function theorem makes
$(u,v)\mapsto(p,v)=(f_u(u,v),v)$ a local smooth coordinate map.
Restrict its inverse $(u(p,v),v)$ to an open product of boxes and put
$g(p,v)=f(u(p,v),v)-p^Tu(p,v)$. The chain rule gives
\[
 g_p=-u,\qquad g_v=f_v,\qquad
 g_{vv}=f_{vv}-f_{vu}f_{uu}^{-1}f_{uv}=0.
\]
The last equality follows from constant rank and the Schur complement.
Thus, on the product chart,
\begin{equation}\label{eq:legendre-fiber}
 g(p,v)=b(p)^Tv+c(p),\qquad
 u(p,v)=-Db(p)^Tv-\nabla c(p).
\end{equation}
For fixed $p$ the fiber is affine in $v$, its gradient is
$(p,b(p))$, and the affine function
$\ell_p(u,v)=p^Tu+b(p)^Tv+c(p)$ agrees with both $f$ and its gradient
along that fiber. Convexity of $f$ is unnecessary for this differential
use of Legendre coordinates.

Choose smaller closed parameter boxes $P_0,V_0$ within the chart and
$P_1$ with $P_0\subset\operatorname{int}P_1$, still within it.
There are constants $L,\eta,M>0$ such that $u(p,v)$ is $L$-Lipschitz
in $p$ in the maximum norm on $P_1\times V_0$, and a transverse
$\eta$-neighborhood of this compact image lies in $U$ with Hessian
operator norm at most $M$. For small $h$, cover $P_0$ by
$O(h^{-r})$ cubes of radius $h$ and centers $p_\nu\in P_1$. Form
\[
 C_\nu(h)=\{(u,v):v\in V_0,\quad
                  \|u-u(p_\nu,v)\|_\infty\le Lh\}.
\]
This is a bounded polytope by \eqref{eq:legendre-fiber}, and the
Lipschitz bound shows that the polytopes cover the image of
$P_0\times V_0$. For any point in a tube compare it with the central
fiber point of the same $v$. Function and gradient errors vanish there,
and the displacement has only $r$ coordinates, each at most $Lh$.
Taylor's theorem gives, throughout the \emph{whole} tube,
\[
 |f-\ell_{p_\nu}|\le\tfrac12MrL^2h^2,
\]
provided $Lh\le\eta$. This estimate does not assume that the tube
point's own gradient parameter lies in the selected parameter cube.

The interiors of images of smaller parameter products cover $B$.
Compactness selects finitely many charts, including boundary points
because $U$ is a neighborhood. Intersect all tubes with $B$ and use
uniform constants across this finite collection. There are
$O(h^{-r})$ bounded polytopes covering $B$. Choose $h$ proportional
to $\sqrt\eps$ so the Taylor error is at most $\eps/2$, and impose
$|w-\ell_{p_\nu}(x)|\le\eps/2$ on each. Their union contains the graph
and has error at most $\eps$. Apply \cref{lem:disjunction}.
\end{proof}

For $f(x)=\|x\|_2$ on $[1,2]^n$, $n\ge2$, the Hessian is
$\|x\|^{-1}(I-xx^T/\|x\|^2)$ and has constant rank $n-1$.
Thus its coefficient is $(n-1)/2$ although its null line varies.
In dimension two the global Hessian span is the whole symmetric
matrix space: three distinct ratios $x/y$ give independent vectors
of coefficients of $(y^2,-xy,x^2)$ by the Vandermonde determinant.
Thus $r_{\mathrm{all}}=2$ while the exact coefficient is $1/2$.
An elementary upper construction uses angular sectors of width $h$.
On a sector with central unit direction $u$ and bounded radius $R$,
$0\le\|x\|-u^Tx\le R(1-\cos(h/2))\le Rh^2/8$.
Intersecting the sectors with the positive box gives polyhedral cells,
and their affine bands use $O(\eps^{-1/2})$ members.

The theorem assumes constant scalar rank on a neighborhood of the
whole box. Rank-changing singularities and general vector maps with
$r_{\mathrm{loc}}<r_{\mathrm{all}}$ are not characterized by these
results. A nonlinear coordinate substitution is not a free operation
on convex lifts; the proof uses its charts only to select actual
polyhedra in the original coordinates.

\subsection{Positive perspectives}\label{subsec:perspective}

\begin{proposition}[Perspective transfer]\label{prop:perspective}
Let $F:D\to\R^m$ have a binary linear graph relaxation of component
error $\delta$. For $0<\ell\le u<\infty$ define
\[
 G(x,t)=tF(x/t),\qquad
 D_{\mathrm{pers}}=\{(x,t):\ell\le t\le u,\ x/t\in D\}.
\]
There is a binary linear graph relaxation of $G$ with the same binaries
and error at most $u\delta$. For each $t_0\in[\ell,u]$,
\begin{align*}
 \pconv(F,D;\eps/t_0)&\le\pconv(G,D_{\mathrm{pers}};\eps),\\
 \pbin(F,D;\eps/t_0)&\le\pbin(G,D_{\mathrm{pers}};\eps)
                         \le\pbin(F,D;\eps/u),
\end{align*}
where scalar tolerances denote componentwise boxes. Hence a matching
logarithmic coefficient for $F$ is preserved by $G$.
\end{proposition}
\begin{proof}
Write the original rows, including the domain restrictions, as
$Ay+Bv+Ca+Ez\le h$, with $z\in\{0,1\}^p$.
Introduce $a'$ and $v'=tz$, imposing the exact binary-product rows
\[
 \ell z_i\le v'_i\le uz_i,\qquad
 t-u(1-z_i)\le v'_i\le t-\ell(1-z_i).
\]
Replace the original rows by $Ax+Bw+Ca'+Ev'\le ht$ and impose
$\ell\le t\le u$. Division by positive $t$ recovers the original
formulation with $y=x/t,v=w/t,a=a'/t$, and conversely every original
feasible point scales to one of these points. Thus graph containment
is exact and $|w_j-tF_j(x/t)|\le t\delta\le u\delta$.
The transformation adds $p$ continuous product variables and $4p+2$
rows, besides the new input $t$; original continuous auxiliaries need
no finite bounds. Rational data retain polynomial encoding size.
Finally, fixing $t=t_0$ and dividing visible variables by $t_0$ gives
the two lower transfers for arbitrary convex or binary linear lifts.
\end{proof}

Perspective geometry is classical \citep[Sections 2.3.3 and 3.2.6]{boyd2004}.
Here the relevant fact is that homogenizing rows adds no binary
coordinates. The product encoding is not valid for an unrestricted
integer variable in place of a binary, and no such upper transfer is
asserted. The positive lower bound on $t$ excludes the cone apex.

For quadratic $F$, the perspective outputs are
\[
 G_j(x,t)=\frac{x^TH_jx}{2t}+a_j^Tx+b_jt.
\]
By \cref{thm:ncrank,prop:perspective}, their precision coefficient is
$r/2$ and their binary upper formulation has size
$O(1+\log(1/\eps))$ for fixed data. The domain is the bounded polytope
$\ell\le t\le u$, $l_it\le x_i\le u_it$ when
$D=\prod_i[l_i,u_i]$. The same coefficient holds on an ordinary
full-dimensional $(x,t)$ box with positive $t$: enclose all ratios in
a larger fixed box for the upper bound, retain the requested input
box, and fix $t=t_0$ for the lower bound.
In particular $x^2/t$ on $[0,1]\times[1,2]$ has coefficient $1/2$.
Its Hessian is $(2/t^3)(t,-x)(t,-x)^T$, of rank one; two nonparallel
such vectors give a positive definite sum, so the global Hessian
span has rank two. This provides a rational, compactly formulable
example in which the global smooth-span upper invariant is strictly
larger than the exact coefficient.
